\begin{notechapter}
\chapter{Uniform Convergence and the Symmetrization Lemma}\label{ch:symmetrization-lemma}
\chaptersubtitle{From finite to infinite hypothesis classes}

\begin{chapterintro}
\Cref{ch:erm} reduced the success of empirical risk minimization to a single
question: is the empirical risk close to the risk \emph{uniformly} over the
hypothesis class? For a finite class, Hoeffding's inequality and a union bound
answered it. This chapter answers it for infinite classes, such as all linear
classifiers. You will learn
\begin{itemize}[nosep]
  \item why the obvious union-bound argument fails;
  \item the \emph{symmetrization lemma}, which replaces the unknown risk by the
        empirical risk on an independent ``ghost'' sample, with a complete
        proof;
  \item how random signs (Rademacher variables) turn the problem into one about
        finitely many patterns;
  \item the \emph{shattering coefficient} and the \emph{VC dimension}, with
        examples;
  \item the Vapnik--Chervonenkis generalization bound, and what it says in
        numbers.
\end{itemize}
\prerequisites{\cref{ch:erm} (in particular \cref{erm:thm:uniform} and
Hoeffding's inequality); conditional expectation; Chebyshev's inequality.}
\end{chapterintro}

\section{The plan}\label{symmetrization-lemma:sec:plan}

For a finite class we bounded $\Prob(\sup_f \abs{\ERisk(f) - \Risk(f)} \geq \epsilon)$
by summing over the $\abs{\Hs}$ hypotheses. With infinitely many hypotheses the
sum is infinite. Yet on $n$ data points a classifier is only seen through its
$n$ predictions, and even an infinite class of binary classifiers produces at
most $2^n$ different prediction patterns. The difficulty is that the risk
$\Risk(f)$ is \emph{not} determined by the predictions on the sample: it
depends on $f$ everywhere.

The argument therefore has two steps.
\begin{enumerate}
  \item \textbf{Symmetrization.} Replace the risk $\Risk(f)$ by the empirical
        risk on a second, independent sample of the same size (the
        \emph{ghost sample}). Afterwards only the values of $f$ on $2n$ points
        matter.
  \item \textbf{Counting.} On $2n$ points the class produces finitely many
        patterns; apply Hoeffding's inequality to each and a union bound over
        the patterns.
\end{enumerate}
The ghost sample is a device of the proof: it is never collected.

\section{Setting and notation}\label{symmetrization-lemma:sec:notation}

Write $Z = (X, Y)$ for a data point and $Z_1,\dots,Z_n$ for the i.i.d.\ sample.
For each hypothesis $f$, the function $g(z) = \ell(f(x), y)$ is its loss; let
$\mathcal G = \{g_f : f\in\Hs\}$ be the class of losses, which we assume take
values in $[0,1]$ (the 0--1 loss is the main example). Following the notation
of empirical process theory,\footnote{Other texts write $P^{\text{true}}g$,
$P^{\text{emp}}g$ and $P^{\text{emp}'}g$ for these three quantities.}
\[
  Pg = \E[g(Z)] = \Risk(f), \qquad
  P_n g = \frac1n\sum_{i=1}^n g(Z_i) = \ERisk(f), \qquad
  P'_n g = \frac1n\sum_{i=1}^n g(Z'_i),
\]
where the ghost sample $Z'_1, \dots, Z'_n$ is i.i.d.\ with the same distribution
as the $Z_i$, and independent of them. For a fixed $g$,
$\E[P'_n g] = \E[P_n g] = Pg$. We assume the suprema below are measurable
(true, for instance, when $\mathcal G$ is countable or can be approximated by a
countable subclass).

\section{Why the obvious argument fails}\label{symmetrization-lemma:sec:naive}

The tempting route is the triangle inequality. For every $g$,
\[
  \abs{P_n g - Pg} \leq \abs{P_n g - P'_n g} + \abs{P'_n g - Pg},
\]
so with the events
\[
  A = \Bigl\{\sup_g \abs{P_n g - Pg} > \epsilon\Bigr\},\quad
  B = \Bigl\{\sup_g \abs{P_n g - P'_n g} > \tfrac{\epsilon}{2}\Bigr\},\quad
  C = \Bigl\{\sup_g \abs{P'_n g - Pg} > \tfrac{\epsilon}{2}\Bigr\},
\]
we have $A \subseteq B \cup C$: if neither $B$ nor $C$ occurs, every
$\abs{P_ng - Pg}$ is at most $\frac\epsilon2 + \frac\epsilon2$. The union bound
gives $\Prob(A) \leq \Prob(B) + \Prob(C)$.

\begin{pitfall}[This inequality is true but useless]
The ghost sample has the same distribution as the sample, so
$\Prob(C) = \Prob\bigl(\sup_g\abs{P_n g - Pg} > \frac\epsilon2\bigr)$, which is
\emph{at least} $\Prob(A)$ because its threshold is smaller. The bound
$\Prob(A) \leq \Prob(B) + \Prob(C)$ is therefore weaker than the trivial
$\Prob(A)\leq\Prob(C)$, and one cannot deduce $\Prob(A) \leq 2\Prob(B)$ from it.
The union bound treats the two deviations as if they were unrelated. The real
proof uses the \emph{independence} of the ghost sample: once the first sample
is fixed, the ghost sample is close to the truth with high probability.
\end{pitfall}

\section{The symmetrization lemma}\label{symmetrization-lemma:sec:lemma}

\begin{lemma}[Symmetrization]\label{symmetrization-lemma:lem:main}
Let the functions of $\mathcal G$ take values in $[0,1]$. For every $t > 0$
with $nt^2 \geq 2$,
\[
  \Prob\Bigl[\sup_{g\in\mathcal G}\,(Pg - P_n g) \geq t\Bigr]
  \;\leq\; 2\,\Prob\Bigl[\sup_{g\in\mathcal G}\,(P'_n g - P_n g) \geq \tfrac t2\Bigr].
\]
The left-hand probability is over the sample; the right-hand one over the
sample and the ghost sample.
\end{lemma}

\begin{intuition}[Bad for one sample, good for the other]
Suppose the sample is ``bad'': some $g$ has an empirical risk at least $t$ below
its true risk. That $g$ was selected using the sample only, so it is a fixed
function from the point of view of the ghost sample, and for a fixed function
the ghost empirical risk is close to the true risk with probability at least
$\frac12$. Whenever that happens, the two samples disagree about $g$ by at
least $\frac t2$. So a bad sample leads, at least half of the time, to a visible
disagreement between two samples --- an event that involves no unknown
quantity. The factor $2$ in the lemma is $1/\frac12$.
\end{intuition}

\begin{proof}
Let $A = \{\sup_g (Pg - P_n g) \geq t\}$ be the event on the left.

\emph{Step 1: a near-maximizer.} The supremum need not be attained. Fix
$\epsilon\in(0,t)$. For each sample, choose $g_\epsilon \in \mathcal G$ with
\[
  Pg_\epsilon - P_n g_\epsilon \geq \sup_{g\in\mathcal G}(Pg - P_ng) - \epsilon .
\]
The function $g_\epsilon$ depends on $Z_1,\dots,Z_n$ but not on the ghost
sample. Let $A_\epsilon = \{Pg_\epsilon - P_n g_\epsilon \geq t-\epsilon\}$. If
$A$ occurs, the supremum is at least $t$, so $g_\epsilon$ achieves at least
$t - \epsilon$: $A\subseteq A_\epsilon$ and
\begin{equation}\label{symmetrization-lemma:eq:step1}
  \Prob(A) \leq \Prob(A_\epsilon).
\end{equation}

\emph{Step 2: a deterministic inclusion.} Let $s = \frac{t-\epsilon}{2}$. If
$A_\epsilon$ occurs and the ghost sample is good for $g_\epsilon$, in the sense
that $Pg_\epsilon - P'_n g_\epsilon < s$, then
\[
  P'_n g_\epsilon - P_n g_\epsilon
  = \underbrace{(Pg_\epsilon - P_n g_\epsilon)}_{\geq\, 2s}
  - \underbrace{(Pg_\epsilon - P'_n g_\epsilon)}_{<\, s} > s .
\]
In terms of indicators, for every pair of samples,
\begin{equation}\label{symmetrization-lemma:eq:indicators}
  \bbone_{A_\epsilon}\;\bbone\bigl[Pg_\epsilon - P'_n g_\epsilon < s\bigr]
  \;\leq\; \bbone\bigl[P'_n g_\epsilon - P_n g_\epsilon \geq s\bigr].
\end{equation}

\emph{Step 3: average over the ghost sample.} Condition on $Z_1,\dots,Z_n$.
Then $\bbone_{A_\epsilon}$ is a constant, $g_\epsilon$ is a fixed function,
and the ghost sample is still an i.i.d.\ sample independent of everything
else. Taking the conditional expectation of \eqref{symmetrization-lemma:eq:indicators}
(the expectation of an indicator is a probability),
\begin{equation}\label{symmetrization-lemma:eq:conditional}
  \bbone_{A_\epsilon}\;\Prob'\bigl[Pg_\epsilon - P'_n g_\epsilon < s\bigr]
  \leq \Prob'\bigl[P'_n g_\epsilon - P_n g_\epsilon \geq s\bigr],
\end{equation}
where $\Prob'$ is the probability over the ghost sample with the first sample
fixed. Since $g_\epsilon$ is fixed and $g_\epsilon(Z'_i)\in[0,1]$, the
one-sided Hoeffding inequality bounds the lower tail of the ghost empirical
mean:
\[
  \Prob'\bigl[Pg_\epsilon - P'_n g_\epsilon \geq s\bigr] \leq e^{-2ns^2}
  = e^{-n(t-\epsilon)^2/2},
  \quad\text{so}\quad
  \Prob'\bigl[Pg_\epsilon - P'_n g_\epsilon < s\bigr] \geq 1 - e^{-n(t-\epsilon)^2/2}.
\]
Plugging this into \eqref{symmetrization-lemma:eq:conditional} gives, for
every first sample,
\[
  \bbone_{A_\epsilon}\Bigl(1 - e^{-n(t-\epsilon)^2/2}\Bigr)
  \leq \Prob'\bigl[P'_n g_\epsilon - P_n g_\epsilon \geq s\bigr].
\]

\emph{Step 4: average over the first sample.} Taking the expectation over
$Z_1,\dots,Z_n$, the left side becomes
$\Prob(A_\epsilon)(1 - e^{-n(t-\epsilon)^2/2})$ and the right side the
probability, over both samples, that $P'_ng_\epsilon - P_ng_\epsilon\geq s$.
Since $g_\epsilon\in\mathcal G$, that event is contained in
$\{\sup_g (P'_n g - P_n g) \geq s\}$. Combined with
\eqref{symmetrization-lemma:eq:step1},
\begin{equation}\label{symmetrization-lemma:eq:step4}
  \Prob(A)\Bigl(1 - e^{-n(t-\epsilon)^2/2}\Bigr)
  \leq \Prob\Bigl[\sup_g (P'_n g - P_n g) \geq \tfrac{t-\epsilon}{2}\Bigr]
  \qquad\text{for every } \epsilon\in(0,t).
\end{equation}

\emph{Step 5: let $\epsilon\downarrow 0$.} On the left,
$1 - e^{-n(t-\epsilon)^2/2} \to 1 - e^{-nt^2/2}$, and $nt^2\geq2$ gives
$e^{-nt^2/2}\leq e^{-1} < \frac12$. On the right, the events
$\{\sup_g(P'_ng - P_ng) \geq \frac{t-\epsilon}{2}\}$ shrink as $\epsilon$
decreases and their intersection is $\{\sup_g(P'_ng - P_ng) \geq \frac t2\}$;
by continuity of probability from above, their probabilities converge to the
probability of that event. Passing to the limit in
\eqref{symmetrization-lemma:eq:step4},
\[
  \tfrac12\,\Prob(A) \leq \Bigl(1 - e^{-nt^2/2}\Bigr)\Prob(A)
  \leq \Prob\Bigl[\sup_g (P'_n g - P_n g) \geq \tfrac t2\Bigr]. \qedhere
\]
\end{proof}

\begin{pitfall}[No limit of the events $A_\epsilon$ is needed]
The near-maximizer $g_\epsilon$ changes with $\epsilon$, so the events
$A_\epsilon$ need not be monotone nor converge to anything, and one should not
try to prove $\Prob(A_\epsilon)\to\Prob(A)$. The proof never needs it: the
inclusion $A\subseteq A_\epsilon$ is enough, and the only limit concerns the
right-hand side, where the events \emph{are} monotone.
\Cref{ch:convergence-of-events} discusses limits of events in detail.
\end{pitfall}

\begin{remark}[Where $nt^2\geq2$ comes from]
The classical proof uses Chebyshev's inequality instead of Hoeffding's (and
assumes for simplicity that the supremum is attained, so that $\epsilon=0$).
A variable with values in $[0,1]$ has variance at most $\frac14$, so
$\Var'(P'_n g) \leq \frac1{4n}$ and
\[
  \Prob'\bigl[Pg - P'_ng \geq \tfrac t2\bigr] \leq \frac{\Var'(P'_n g)}{(t/2)^2} \leq \frac{1}{nt^2}.
\]
This is at most $\frac12$ exactly when $nt^2\geq2$; Hoeffding's inequality gives
the same conclusion under the same condition. Without some such condition the
lemma is false: for very small $n$ the ghost sample is not reliably good.
\end{remark}

\subsection{The two-sided version}

The same argument, with absolute values throughout, gives the version used for
generalization bounds.

\begin{corollary}[Two-sided symmetrization]\label{symmetrization-lemma:cor:two-sided}
Under the same assumptions, for every $\epsilon>0$ with $n\epsilon^2\geq2$,
\[
  \Prob\Bigl[\sup_{g}\abs{P_n g - Pg} > \epsilon\Bigr]
  \leq 2\,\Prob\Bigl[\sup_{g}\abs{P_n g - P'_n g} > \tfrac\epsilon2\Bigr].
\]
\end{corollary}

\begin{proof}[Proof sketch]
Repeat the proof with $\abs{Pg - P_n g}$ in place of $Pg - P_ng$ and the strict
thresholds of the statement. In Step~3 one now needs
$\Prob'(\abs{P'_n g - Pg} \geq \frac\epsilon2) \leq \frac12$ for a fixed $g$:
the two-sided Chebyshev bound of the remark gives exactly $\frac1{n\epsilon^2}\leq\frac12$.
(The two-sided Hoeffding bound $2e^{-n\epsilon^2/2}$ would require the
stronger condition $n\epsilon^2\geq 2\log 4$.)
\end{proof}

\section{Random signs}\label{symmetrization-lemma:sec:rademacher}

After symmetrization, the quantity to control is
$\sup_g \abs{\frac1n\sum_i (g(Z_i) - g(Z'_i))}$. A second trick makes it
tractable.

Let $\sigma_1,\dots,\sigma_n$ be independent \emph{Rademacher} variables,
equal to $+1$ or $-1$ with probability $\frac12$ each, independent of the
data. Exchanging $Z_i$ and $Z'_i$ does not change the joint distribution of the
two samples, because the pairs are i.i.d.\ and the two samples play symmetric
roles. The exchange flips the sign of $g(Z_i) - g(Z'_i)$ \emph{simultaneously
for all} $g$. Exchanging each pair independently with probability $\frac12$
therefore shows that the two random processes
\[
  \Bigl(\frac1n\sum_{i=1}^n \bigl(g(Z_i) - g(Z'_i)\bigr)\Bigr)_{g\in\mathcal G}
  \quad\text{and}\quad
  \Bigl(\frac1n\sum_{i=1}^n \sigma_i\bigl(g(Z_i) - g(Z'_i)\bigr)\Bigr)_{g\in\mathcal G}
\]
have the same distribution, hence so do their suprema.

\begin{pitfall}[An identity, not an inequality]
Inserting the signs is an exact identity in distribution. It is not a
``cancellation of bias'': the differences $g(Z_i) - g(Z'_i)$ already have mean
zero. The benefit is elsewhere: conditionally on the $2n$ data points, the only
randomness left is in the signs, which are independent and bounded, so
Hoeffding's inequality applies to $\sum_i\sigma_i a_i$ for any fixed numbers
$a_i$.
\end{pitfall}

\section{Counting patterns: shattering and VC dimension}\label{symmetrization-lemma:sec:vc}

Let $\Hs$ be a class of binary classifiers $f:\Xs\to\{0,1\}$.

\begin{definition}[Shattering coefficient and VC dimension]
The \emph{shattering coefficient} (or growth function) of $\Hs$ is the largest
number of distinct labellings that $\Hs$ produces on $n$ points:
\[
  S(\Hs, n) = \max_{x_1,\dots,x_n\in\Xs}\;\Bigl|\bigl\{(f(x_1),\dots,f(x_n)) : f\in\Hs\bigr\}\Bigr| \;\leq\; 2^n .
\]
A set of $n$ points is \emph{shattered} if $\Hs$ realizes all $2^n$ labellings
on it. The \emph{VC dimension} $d_{\mathrm{VC}}(\Hs)$ is the largest $n$ such
that some set of $n$ points is shattered, i.e.\ such that $S(\Hs,n) = 2^n$
(it is $\infty$ if there is no largest such $n$).
\end{definition}

To show $d_{\mathrm{VC}} = d$ one must exhibit \emph{one} set of $d$ points that
is shattered, and show that \emph{no} set of $d+1$ points is.

\begin{example}[Thresholds on the line]
$\Hs = \{x\mapsto\bbone[x\geq\theta] : \theta\in\R\}$. On sorted points
$x_1<\dots<x_n$ a threshold labels a prefix $0$ and the rest $1$, so the
possible labellings are determined by the position of $\theta$ among the
points: $S(\Hs,n) = n+1$. One point is shattered, but on two points
$x_1 < x_2$ the labelling $(1,0)$ is impossible: $d_{\mathrm{VC}} = 1$.
\end{example}

\begin{example}[Intervals]
$\Hs = \{x\mapsto\bbone[a\leq x\leq b]\}$. A labelling is realizable iff its
$1$s are consecutive. Counting the empty labelling and the choices of a first
and last $1$, $S(\Hs,n) = 1 + n + \binom n2$. Two points are shattered; on three
points the pattern $(1,0,1)$ is impossible: $d_{\mathrm{VC}} = 2$.
\end{example}

\begin{example}[Half-planes]\label{symmetrization-lemma:ex:halfplanes}
$\Hs = \{x\mapsto\bbone[w^\top x + b\geq0] : w\in\R^2, b\in\R\}$. Three points
not on a line can be shattered: any subset of them can be cut off from the
others by a line. Four points cannot. If one of them lies inside the triangle
formed by the other three, the labelling ``1 on the triangle, 0 inside'' is
impossible, because a half-plane is convex. Otherwise the four points form a
convex quadrilateral, and labelling the two diagonals differently (the ``XOR''
pattern of \cref{symmetrization-lemma:fig:xor}) is impossible: the two
segments joining same-label points cross, and their intersection point would
have to receive both labels. (If three of the points are on a line, the
middle one cannot be labelled differently from the two others.) Hence
$d_{\mathrm{VC}} = 3$. More generally, affine
classifiers on $\R^d$ have VC dimension $d+1$.
\end{example}

\begin{figure}[ht]
\centering
\begin{tikzpicture}[scale=1.1]
  \draw[gray!60] (0,0) -- (2,2) (0,2) -- (2,0);
  \foreach \p in {(0,0),(2,2)} \fill[bookaccent] \p circle (3.5pt);
  \foreach \p in {(0,2),(2,0)} \draw[thick, bookpitfall, fill=white] \p circle (3.5pt);
  \node[right, font=\small, align=left] at (2.6,1) {filled: label 1\\ open: label 0};
\end{tikzpicture}
\caption{The XOR labelling of four points in convex position cannot be produced
by a half-plane.}
\label{symmetrization-lemma:fig:xor}
\end{figure}

The decisive fact is that the growth function is either \emph{exponential}
forever or eventually \emph{polynomial}.

\begin{lemma}[Sauer--Shelah]\label{symmetrization-lemma:lem:sauer}
If $d_{\mathrm{VC}}(\Hs) = d < \infty$, then for all $n$,
\[
  S(\Hs, n) \leq \sum_{k=0}^d\binom nk, \qquad\text{and}\qquad S(\Hs,n)\leq\Bigl(\frac{en}{d}\Bigr)^d \text{ for } n\geq d .
\]
\end{lemma}

The intervals example attains the first bound with equality. (A proof by
induction on $n$ can be found in any learning theory textbook, e.g.\
Shalev-Shwartz and Ben-David, \emph{Understanding Machine Learning}, Lemma~6.10.)

\section{The VC generalization bound}\label{symmetrization-lemma:sec:bound}

We can now assemble the pieces.

\begin{theorem}[Vapnik--Chervonenkis]\label{symmetrization-lemma:thm:vc}
Let $\Hs$ be a class of binary classifiers, with the 0--1 loss. For every
$\epsilon>0$ with $n\epsilon^2\geq2$,
\[
  \Prob\Bigl(\sup_{f\in\Hs}\abs{\ERisk(f) - \Risk(f)} > \epsilon\Bigr)
  \leq 4\,S(\Hs,2n)\,e^{-n\epsilon^2/8}.
\]
\end{theorem}

\begin{proof}
\emph{Symmetrize.} By \cref{symmetrization-lemma:cor:two-sided}, the
probability is at most
$2\,\Prob\bigl(\sup_g\abs{\frac1n\sum_i(g(Z_i) - g(Z'_i))} > \frac\epsilon2\bigr)$.

\emph{Insert random signs.} By \cref{symmetrization-lemma:sec:rademacher}
this equals
$2\,\Prob\bigl(\sup_g\abs{\frac1n\sum_i\sigma_i(g(Z_i) - g(Z'_i))} > \frac\epsilon2\bigr)$.

\emph{Condition on the data and count.} Fix the $2n$ points
$Z_1,\dots,Z_n,Z'_1,\dots,Z'_n$. The 0--1 loss $g_f(z) = \bbone[f(x)\neq y]$ on
these points is determined by the labelling that $f$ gives to the $2n$ inputs,
so as $f$ ranges over $\Hs$, the vector
$a = \bigl(g(Z_i) - g(Z'_i)\bigr)_{i\leq n}$ takes at most $S(\Hs,2n)$ distinct
values. For one fixed vector $a\in[-1,1]^n$, the terms $\sigma_ia_i$ are
independent, have mean zero and lie in $[-1,1]$, so Hoeffding's inequality
(range $b - a = 2$) gives
\[
  \Prob_\sigma\Bigl(\Bigl|\frac1n\sum_i\sigma_ia_i\Bigr| > \frac\epsilon2\Bigr)
  \leq 2\exp\Bigl(-\frac{2n(\epsilon/2)^2}{4}\Bigr) = 2e^{-n\epsilon^2/8}.
\]
A union bound over the at most $S(\Hs,2n)$ vectors bounds the conditional
probability of the supremum by $2S(\Hs,2n)e^{-n\epsilon^2/8}$. This bound does
not depend on the data, so it also bounds the unconditional probability.
Multiplying by the factor 2 of the first step gives the result.
\end{proof}

Inverting the bound and using Sauer's lemma gives the familiar form.

\begin{corollary}[Generalization bound]\label{symmetrization-lemma:cor:bound}
Let $d = d_{\mathrm{VC}}(\Hs)<\infty$ and $\delta\in(0,1)$. With probability at
least $1-\delta$, simultaneously for all $f\in\Hs$,
\[
  \Risk(f) \leq \ERisk(f) + \sqrt{\frac8n\Bigl(d\log\frac{2en}{d} + \log\frac4\delta\Bigr)},
\]
and the ERM solution satisfies $\Risk(\hat f_n) - \Risk(f_\Hs) \leq 2\epsilon_n$, where
$\epsilon_n$ is the square root above.
\end{corollary}

\begin{proof}
Set $4S(\Hs,2n)e^{-n\epsilon^2/8} = \delta$ and solve:
$\epsilon^2 = \frac8n(\log S(\Hs,2n) + \log\frac4\delta)$. Sauer's lemma gives
$\log S(\Hs,2n) \leq d\log(2en/d)$ (for $2n \geq d$). The condition
$n\epsilon^2\geq 2$ holds automatically, since $n\epsilon^2\geq 8\log 4 > 2$.
The ERM statement follows from \cref{erm:thm:uniform}.
\end{proof}

\begin{example}[The bound in numbers]
For half-planes in the plane ($d = 3$) and $\delta = 0.05$, the deviation
$\epsilon_n$ is about $0.46$ for $n = 1000$, $0.16$ for $n = 10\,000$ and
$0.057$ for $n = 100\,000$. The constants are pessimistic --- in practice the
gap between training and test error of a linear classifier is far smaller ---
but the \emph{rate} is the message: the deviation decreases like
$\sqrt{d\log n/n}$, and the class enters only through its VC dimension, not
through its (infinite) cardinality.
\end{example}

\begin{keyidea}[Learnability]
If $\Hs$ has finite VC dimension, ERM over $\Hs$ is consistent for
\emph{every} distribution $P$, with an explicit rate. Conversely (the
``fundamental theorem of statistical learning''), a class of infinite VC
dimension cannot be learned with guarantees that hold uniformly over all
distributions. The VC dimension is thus the right measure of the size of a
class of classifiers.
\end{keyidea}

\begin{chaptersummary}
\begin{itemize}[nosep]
  \item A naive triangle-inequality and union-bound argument cannot relate
        the deviation from the true risk to the deviation between two samples.
  \item The symmetrization lemma does: a sample on which some hypothesis
        looks too good leads, with probability at least $\frac12$, to a
        disagreement of size $\frac t2$ with an independent ghost sample. The
        condition $nt^2\geq2$ guarantees that probability.
  \item Random signs can be inserted at no cost (identity in distribution);
        conditionally on the $2n$ points, Hoeffding's inequality applies.
  \item On $2n$ points a class of classifiers produces at most $S(\Hs,2n)$
        patterns; by Sauer's lemma this is polynomial when the VC dimension is
        finite.
  \item Together: $\Prob(\sup_f\abs{\ERisk(f)-\Risk(f)}>\epsilon)\leq4S(\Hs,2n)e^{-n\epsilon^2/8}$,
        a deviation of order $\sqrt{d\log n/n}$.
\end{itemize}
\end{chaptersummary}

\section*{Exercises}
\addcontentsline{toc}{section}{Exercises}

\begin{exercise}
Show that the class of indicator functions of axis-aligned rectangles in
$\R^2$ has VC dimension $4$.
\hint{for the upper bound, among five points consider the leftmost,
rightmost, highest and lowest.}
\end{exercise}

\begin{exercise}
Prove the first bound of \cref{symmetrization-lemma:lem:sauer} for the class
of thresholds and for the class of intervals by direct counting, and check
that it is attained.
\end{exercise}

\begin{exercise}
Where exactly does the proof of \cref{symmetrization-lemma:lem:main} use the
independence of the ghost sample from the first sample? Construct an example
where the ghost sample is a copy of the first sample ($Z'_i = Z_i$) and the
conclusion fails.
\end{exercise}

\begin{exercise}
Using \cref{symmetrization-lemma:cor:bound}, find the smallest $n$ (to within
a factor 2) that guarantees an estimation error below $0.05$ with probability
$0.95$ for half-planes in $\R^2$.
\end{exercise}
\end{notechapter}
